% =====================================================================
%  Interpolación en la sincronización Rössler (esquema de Pecora-Carroll)
%  Análisis matemático del módulo de sincronización del sistema esclavo
%  Archivo fuente analizado: Esclavo_TLS KEYSTREAM.py
% =====================================================================
\documentclass[12pt,a4paper]{article}

\usepackage[utf8]{inputenc}
\usepackage[T1]{fontenc}
% es-noshorthands desactiva los caracteres activos de babel-spanish («<», «>»),
% que de otro modo entran en conflicto con la sintaxis de flechas de TikZ.
\usepackage[spanish,es-tabla,es-noshorthands]{babel}
\usepackage{amsmath}
\usepackage{amsfonts}
\usepackage{amssymb}
\usepackage{graphicx}
\usepackage{float}
\usepackage{booktabs}
\usepackage{geometry}
\usepackage{tikz}
\usepackage[hidelinks]{hyperref}

\geometry{left=2.5cm,right=2.5cm,top=2.5cm,bottom=2.5cm}

% --- Entornos de enunciado (definidos con el núcleo de LaTeX, sin amsthm) ---
% Cuerpo en cursiva: teoremas, proposiciones, lemas.
\newtheorem{teorema}{Teorema}[section]
\newtheorem{proposicion}[teorema]{Proposición}
\newtheorem{lema}[teorema]{Lema}

% Cuerpo en redonda: definiciones y observaciones.
% El envoltorio reenvía el título opcional al entorno subyacente y restituye
% la fuente redonda tras la cabecera.
\newtheorem{definicionaux}[teorema]{Definición}
\newenvironment{definicion}[1][]%
  {\if\relax\detokenize{#1}\relax \definicionaux \else \definicionaux[#1]\fi
   \normalfont}%
  {\enddefinicionaux}

\newtheorem{observacionaux}[teorema]{Observación}
\newenvironment{observacion}[1][]%
  {\if\relax\detokenize{#1}\relax \observacionaux \else \observacionaux[#1]\fi
   \normalfont}%
  {\endobservacionaux}

% Entorno de demostración.
\newenvironment{proof}[1][Demostración]%
  {\par\medskip\noindent\textit{#1.}\space\ignorespaces}%
  {\hfill$\square$\par\medskip}

\newcommand{\R}{\mathbb{R}}
\newcommand{\dnorm}[1]{\left\lVert #1 \right\rVert}
\newcommand{\infnorm}[1]{\dnorm{#1}_{\infty}}

\title{Interpolación en la sincronización caótica de Rössler\\
       bajo el esquema de Pecora y Carroll\\[0.35em]
       \large Análisis matemático del módulo de sincronización del sistema esclavo}
\author{Documentación técnica del sistema de cifrado caótico\\
        \texttt{Esclavo\_TLS KEYSTREAM.py}}
\date{\today}

\begin{document}

\maketitle
\tableofcontents
\newpage

% =====================================================================
\section{Objeto del documento}
% =====================================================================

El sistema esclavo debe reconstruir, en su propia máquina, la misma señal caótica
que el maestro utilizó para cifrar. Para lograrlo integra un oscilador de Rössler
acoplado al del maestro mediante el esquema de sincronización de Pecora y Carroll.
El acoplamiento requiere conocer una señal del maestro \emph{en todo instante de
tiempo}; sin embargo, lo único que viaja por la red es una \emph{lista finita de
muestras} de esa señal.

Este documento desarrolla, de forma matemáticamente explícita, el papel que
desempeña la interpolación dentro de ese proceso: qué objeto matemático construye,
qué hueco lógico cubre, y por qué la calidad de esa reconstrucción determina
directamente la calidad de la sincronización alcanzable. El desarrollo se apoya
únicamente en las ecuaciones efectivamente implementadas en el código del sistema
esclavo y en los valores de configuración que este comparte con el maestro.

La línea argumental es la siguiente:

\begin{enumerate}
  \item El esquema de Pecora y Carroll garantiza sincronización \emph{siempre que
        la señal guía esté disponible como función continua del tiempo}
        (Secciones \ref{sec:planteamiento} y \ref{sec:error}).
  \item El canal de comunicación solo entrega muestras discretas, y el integrador
        numérico exige evaluaciones fuera de esa malla
        (Sección \ref{sec:discreto}).
  \item La interpolación es el operador que convierte las muestras en la función
        continua que el esquema necesita, e introduce un error de reconstrucción
        $\delta(t)$ (Sección \ref{sec:interpolacion}).
  \item El error de sincronización queda acotado por una cantidad
        \emph{proporcional} a $\infnorm{\delta}$: la interpolación no es un detalle
        de implementación, sino el término que fija el piso del error
        (Sección \ref{sec:resultado}).
  \item El orden de aproximación del interpolante controla ese piso, y el spline
        cúbico empleado reduce dicha cota de $\mathcal{O}(h)$ a $\mathcal{O}(h^4)$
        (Sección \ref{sec:ordenes}).
  \item Esa reducción se traduce, en la etapa de descifrado, en un margen
        cuantificable frente al umbral de cuantización de los píxeles
        (Sección \ref{sec:cripto}).
\end{enumerate}

Las mediciones experimentales asociadas (errores de sincronización, correlaciones,
distancias de Hamming y tiempos de proceso) se encuentran documentadas por separado
y no se reproducen aquí; este texto se ocupa exclusivamente de la justificación
matemática del mecanismo.

% =====================================================================
\section{Planteamiento: maestro, esclavo y el esquema de Pecora--Carroll}
\label{sec:planteamiento}
% =====================================================================

\subsection{El sistema maestro}

El maestro integra el oscilador de Rössler en su forma autónoma clásica. Escribiendo
el estado como $X_m(t) = \bigl(x_m(t),\,y_m(t),\,z_m(t)\bigr) \in \R^3$:

\begin{equation}
\label{eq:maestro}
\left\{
\begin{aligned}
  \dot{x}_m &= -\,y_m - z_m,\\
  \dot{y}_m &= x_m + a\,y_m,\\
  \dot{z}_m &= b + z_m\,(x_m - c),
\end{aligned}
\right.
\qquad X_m(0) = X_m^{0}.
\end{equation}

Los parámetros $(a,b,c)$ y la condición inicial $X_m^0$ constituyen parte de la
clave del sistema. El campo vectorial de \eqref{eq:maestro} es polinomial, por lo
tanto de clase $C^{\infty}$; en consecuencia toda solución $t \mapsto X_m(t)$ es
también $C^{\infty}$ en su intervalo de existencia. Este hecho, aparentemente
técnico, será decisivo más adelante: es exactamente la hipótesis de regularidad que
permite justificar el uso de un interpolante de alto orden
(Sección \ref{sec:ordenes}).

El maestro utiliza la componente $x_m$ para enmascarar la información, y transmite
la componente $y_m$ como \emph{señal guía}.

\subsection{La descomposición de Pecora y Carroll}

La idea de Pecora y Carroll consiste en partir el sistema en dos bloques: un
\emph{subsistema conductor}, cuya señal se transmite, y un \emph{subsistema
respuesta}, que se replica en el receptor y se deja gobernar por la señal recibida.
En la formulación original, la réplica no integra la ecuación de la variable
conductora: simplemente la sustituye por el valor recibido.

Aquí la partición es
\[
  \underbrace{y}_{\text{conductor}}
  \qquad\Big|\qquad
  \underbrace{(x,\,z)}_{\text{respuesta}} ,
\]
y el subsistema respuesta impulsado por la señal $y_m(t)$ se escribe
\begin{equation}
\label{eq:respuesta}
\left\{
\begin{aligned}
  \dot{x}_s &= -\,y_m(t) - z_s,\\
  \dot{z}_s &= b + z_s\,(x_s - c).
\end{aligned}
\right.
\end{equation}

El criterio de Pecora y Carroll establece que la réplica \eqref{eq:respuesta}
converge a la trayectoria del maestro, cualquiera que sea su condición inicial,
si y solo si todos los \emph{exponentes de Lyapunov condicionales} del subsistema
respuesta son negativos. Estos exponentes miden la tasa de crecimiento o
decaimiento de perturbaciones \emph{transversales} a la variedad de sincronización
\[
  \mathcal{M} \;=\; \bigl\{\,(X_m, X_s) \in \R^3\times\R^3 \;:\; X_s = X_m \,\bigr\},
\]
y son «condicionales» porque se evalúan a lo largo de la trayectoria impuesta por
el conductor, no sobre el sistema aislado.

\subsection{La forma acoplada implementada}

El sistema esclavo no realiza la sustitución literal, sino su versión
\emph{acoplada} o difusiva, que es la que aparece en el código. Con
$X_s(t)=(x_s,y_s,z_s)$:

\begin{equation}
\label{eq:esclavo}
\boxed{
\left\{
\begin{aligned}
  \dot{x}_s &= -\,y_s - z_s,\\
  \dot{y}_s &= x_s + a\,y_s \;+\; k\,\bigl(\,y_m(t) - y_s\,\bigr),\\
  \dot{z}_s &= b + z_s\,(x_s - c),
\end{aligned}
\right.}
\qquad X_s(0) = X_s^{0} \neq X_m^{0}.
\end{equation}

El término $k\,(y_m - y_s)$, con ganancia de acoplamiento $k>0$, penaliza la
discrepancia entre la variable $y$ local y la del maestro. Las dos formulaciones
están relacionadas de manera transparente: al despejar
\[
  \frac{1}{k}\,\dot{y}_s \;=\; \frac{1}{k}\bigl(x_s + a y_s\bigr) + \bigl(y_m(t) - y_s\bigr),
\]
y hacer $k \to \infty$, el miembro izquierdo se anula y se obtiene
$y_s \equiv y_m(t)$, es decir, la sustitución de Pecora y Carroll en sentido
estricto. La forma \eqref{eq:esclavo} es entonces una regularización con ganancia
finita del esquema original, y $k$ actúa como el parámetro que interpola entre
«sistema desacoplado» ($k=0$) y «sustitución exacta» ($k\to\infty$).

\begin{observacion}
Obsérvese que $X_s^0 \neq X_m^0$ de manera deliberada: el maestro y el esclavo
arrancan en puntos distintos del espacio de fases. Que dos sistemas caóticos con
condiciones iniciales diferentes terminen recorriendo la misma trayectoria es,
precisamente, el fenómeno que el esquema explota. En ausencia de acoplamiento la
sensibilidad a condiciones iniciales las separaría exponencialmente.
\end{observacion}

\begin{table}[H]
\centering
\caption{Correspondencia entre la notación matemática y las variables del código.}
\label{tab:notacion}
\begin{tabular}{@{}lll@{}}
\toprule
\textbf{Símbolo} & \textbf{Significado} & \textbf{Identificador en el código} \\
\midrule
$y_m(t_n)$        & Muestras de la señal guía              & \texttt{y\_maestro} \\
$t_n$             & Malla de instantes de muestreo         & \texttt{t\_maestro}, \texttt{times} \\
$\hat{y}_m(t)$    & Guía reconstruida por interpolación    & \texttt{y\_maestro\_interp} \\
$X_s(t)$          & Estado del esclavo                     & \texttt{sol\_esclavo.y} \\
$k$               & Ganancia de acoplamiento               & \texttt{K} \\
$h$               & Paso nominal de la malla               & \texttt{H} \\
$N$               & Número total de muestras               & \texttt{tiempo\_sinc + keystream} \\
$N_{\mathrm{sinc}}$ & Muestras reservadas al transitorio   & \texttt{tiempo\_sinc} \\
$\alpha$          & Factor de escala de la imagen          & \texttt{IMG\_SCALE} \\
$(a,b,c)$         & Parámetros de Rössler (clave)          & \texttt{ROSSLER\_PARAMS} \\
\bottomrule
\end{tabular}
\end{table}

% =====================================================================
\section{La dinámica del error de sincronización}
\label{sec:error}
% =====================================================================

\subsection{El sistema de error es exactamente lineal}

Defínase el error de sincronización como la diferencia de estados
\[
  e(t) \;=\; X_s(t) - X_m(t) \;=\; \bigl(e_x(t),\,e_y(t),\,e_z(t)\bigr)^{\!\top}.
\]

Restando \eqref{eq:maestro} de \eqref{eq:esclavo} componente a componente, y
suponiendo por el momento que el esclavo dispone de la señal guía exacta $y_m(t)$:

\begin{align}
  \dot{e}_x &= -e_y - e_z, \label{eq:ex}\\[0.2em]
  \dot{e}_y &= e_x + a\,e_y + k\bigl(y_m - y_s\bigr)
             = e_x + (a-k)\,e_y, \label{eq:ey}\\[0.2em]
  \dot{e}_z &= z_s(x_s-c) - z_m(x_m-c). \label{eq:ez0}
\end{align}

La tercera ecuación merece un paso intermedio. Sumando y restando $z_m x_s$:
\begin{align}
  \dot{e}_z
   &= z_s x_s - z_m x_m - c\,e_z \nonumber\\
   &= \bigl(z_s - z_m\bigr)x_s + z_m\bigl(x_s - x_m\bigr) - c\,e_z \nonumber\\
   &= \bigl(x_s(t) - c\bigr)\,e_z \;+\; z_m(t)\,e_x. \label{eq:ez}
\end{align}

Nótese que el término constante $b$ se cancela exactamente. Reuniendo
\eqref{eq:ex}, \eqref{eq:ey} y \eqref{eq:ez}:

\begin{equation}
\label{eq:sistema-error}
\dot{e} \;=\; J(t)\,e,
\qquad
J(t) \;=\;
\begin{pmatrix}
   0        & -1   & -1 \\
   1        & a-k  & 0  \\
   z_m(t)   & 0    & x_s(t)-c
\end{pmatrix}.
\end{equation}

Conviene subrayar un punto: \eqref{eq:sistema-error} \emph{no} es una linealización
ni una aproximación de primer orden. Es una identidad exacta. La no linealidad
original quedó absorbida en los coeficientes $z_m(t)$ y $x_s(t)-c$, que son
funciones conocidas del tiempo una vez fijadas las trayectorias. El sistema de
error es, literalmente, un sistema lineal con coeficientes variables.

\subsection{Condición de sincronización}

Sea $\Phi(t,s)$ la matriz de transición de estados asociada a \eqref{eq:sistema-error},
es decir, la solución matricial de
\[
  \frac{\partial}{\partial t}\Phi(t,s) = J(t)\,\Phi(t,s),
  \qquad \Phi(s,s) = I_3 .
\]
La solución del sistema de error es entonces $e(t)=\Phi(t,0)\,e(0)$.

\begin{definicion}[Sincronización exponencial]
Se dice que el esquema sincroniza exponencialmente si existen constantes
$M \ge 1$ y $\lambda > 0$ tales que
\begin{equation}
\label{eq:cota-transicion}
  \dnorm{\Phi(t,s)} \;\le\; M\,e^{-\lambda\,(t-s)},
  \qquad \text{para todo } t \ge s \ge 0 .
\end{equation}
\end{definicion}

La constante $\lambda$ es la tasa de contracción transversal; su existencia
equivale a que los exponentes de Lyapunov condicionales del subsistema respuesta
sean estrictamente negativos, que es el criterio de Pecora y Carroll. Bajo
\eqref{eq:cota-transicion} se tiene de inmediato
\[
  \dnorm{e(t)} \;\le\; M\,e^{-\lambda t}\,\dnorm{e(0)} \;\xrightarrow[t\to\infty]{}\; 0 ,
\]
esto es, sincronización completa para cualquier condición inicial del esclavo.

\begin{observacion}[El papel de la ganancia $k$]
En $J(t)$ la ganancia aparece únicamente en la entrada $(2,2)$, como $a-k$. Elegir
$k > a$ vuelve negativa esa entrada diagonal y es el mecanismo que estabiliza la
dirección transversal; incrementar $k$ tiende a aumentar $\lambda$ y por lo tanto a
acelerar el transitorio. Se verá en la Sección \ref{sec:resultado} que esta ventaja
tiene una contrapartida en presencia de una guía imperfecta.
\end{observacion}

\begin{observacion}[Interpretación del tiempo de sincronización]
La cota $M e^{-\lambda t}\dnorm{e(0)}$ explica por qué el diseño reserva las
primeras $N_{\mathrm{sinc}}$ muestras como periodo de sincronización antes de
extraer el flujo de clave: son la ventana temporal concedida al transitorio para
decaer. La señal de enmascaramiento se toma únicamente a partir del índice
$N_{\mathrm{sinc}}$, cuando el primer término ya se ha vuelto despreciable.
\end{observacion}

% =====================================================================
\section{El obstáculo real: la señal guía es discreta}
\label{sec:discreto}
% =====================================================================

Toda la Sección \ref{sec:error} descansa sobre una suposición que la implementación
no puede satisfacer literalmente: que el esclavo conoce $y_m(t)$ para todo
$t$ real. Conviene examinar con precisión qué es lo que realmente recibe y qué es
lo que realmente necesita.

\subsection{Lo que el esclavo recibe}

El maestro resuelve \eqref{eq:maestro} sobre una malla uniforme de $N$ instantes
y transmite la evaluación de $y_m$ sobre ella. Formalmente, el esclavo recibe el
conjunto finito de pares
\begin{equation}
\label{eq:datos}
  \mathcal{D} \;=\; \bigl\{\,(t_n,\;\eta_n)\,\bigr\}_{n=0}^{N-1},
  \qquad
  \eta_n = y_m(t_n),
  \qquad
  t_n = n\,h ,
\end{equation}
con $h$ el paso de la malla y $N = N_{\mathrm{sinc}} + N_{\mathrm{key}}$. El
conjunto $\mathcal{D}$ es un objeto \emph{discreto}: contiene información sobre
$y_m$ en $N$ puntos y en ningún otro.

\subsection{Lo que el integrador necesita}

El esclavo integra \eqref{eq:esclavo} con un método Runge--Kutta explícito de
paso adaptativo, de orden $5(4)$ del tipo Dormand--Prince. Un paso de este método,
partiendo de $t$ con tamaño $\tau$, evalúa el campo vectorial en las etapas
\begin{equation}
\label{eq:etapas}
  T_i \;=\; t + c_i\,\tau,
  \qquad
  (c_1,\dots,c_7) \;=\;
  \Bigl(0,\;\tfrac{1}{5},\;\tfrac{3}{10},\;\tfrac{4}{5},\;\tfrac{8}{9},\;1,\;1\Bigr),
\end{equation}
y en cada una de ellas debe calcular
\[
  \dot{y}_s\big|_{T_i} \;=\; x_s + a\,y_s + k\,\bigl(\,y_m(T_i) - y_s\,\bigr),
\]
lo que exige el valor de la guía \emph{en el instante $T_i$}.

Aquí aparece el conflicto de manera nítida. Aun en el caso más favorable de paso
fijo $\tau = h$, los nodos intermedios
\[
  t_n + \tfrac{1}{5}h,\quad
  t_n + \tfrac{3}{10}h,\quad
  t_n + \tfrac{4}{5}h,\quad
  t_n + \tfrac{8}{9}h
\]
\emph{no pertenecen} a la malla $\{t_n\}$: son puntos interiores de los intervalos.
Y el método es de paso adaptativo, de modo que $\tau$ se ajusta continuamente
según el control de error y en general $\tau \neq h$; los instantes solicitados son
entonces esencialmente arbitrarios dentro de $[0, T]$.

\begin{figure}[H]
\centering
\begin{tikzpicture}[x=1.4cm,y=1cm,>=stealth]
  % --- fila superior: lo que el canal entrega ---
  \foreach \i in {0,1,2,3,4}{
     \fill (\i,1.55) circle (2.2pt);
     \draw[dashed,gray!55] (\i,1.45) -- (\i,0.16);
  }
  \node[anchor=west,font=\footnotesize] at (4.85,1.55)
       {muestras $\eta_n$ recibidas};

  % --- fila intermedia: lo que el integrador solicita ---
  \foreach \x in {0.21,0.58,0.97,1.34,1.79,2.22,2.68,3.09,3.53,3.87}{
     \draw[->] (\x,0.95) -- (\x,0.15);
  }
  \node[anchor=west,font=\footnotesize,align=left] at (4.85,0.70)
       {instantes $T_i$ exigidos\\por el integrador};

  % --- eje temporal y nodos de la malla ---
  \draw[->,thick] (-0.35,0) -- (4.55,0) node[right,font=\footnotesize] {$t$};
  \foreach \i/\lab in {0/{t_{n}},1/{t_{n+1}},2/{t_{n+2}},3/{t_{n+3}},4/{t_{n+4}}}{
     \draw[thick] (\i,-0.13) -- (\i,0.13);
     \node[below=2pt,font=\footnotesize] at (\i,-0.13) {$\lab$};
  }
\end{tikzpicture}
\caption{La malla de muestreo y las evaluaciones requeridas por las etapas del
método Runge--Kutta no coinciden. La interpolación es el operador que define
$\hat{y}_m$ sobre todo el intervalo, y no solo sobre los nodos.}
\label{fig:malla}
\end{figure}

\subsection{Formulación del problema}

El problema queda planteado en términos precisos:

\begin{quote}
\emph{Construir, a partir del conjunto finito $\mathcal{D}$ de \eqref{eq:datos},
una función $\hat{y}_m : [0,T] \to \R$ definida en todo punto, que pueda evaluarse
en cualquier instante solicitado por el integrador y que aproxime a $y_m$ tan bien
como sea posible.}
\end{quote}

Es exactamente el problema de interpolación. No se trata, por tanto, de un
refinamiento opcional añadido para «suavizar» la señal: sin un operador de
reconstrucción, el sistema \eqref{eq:esclavo} ni siquiera está definido como
ecuación diferencial, porque su miembro derecho carece de valor en casi todo punto
de $[0,T]$.

% =====================================================================
\section{La interpolación como reconstrucción de la señal guía}
\label{sec:interpolacion}
% =====================================================================

\subsection{El operador de reconstrucción}

Sea $\mathcal{I}$ un operador que a cada conjunto de muestras le asigna una función
continua,
\[
  \mathcal{I} : \mathcal{D} \;\longmapsto\; \hat{y}_m = \mathcal{I}[\mathcal{D}] \in C[0,T].
\]
El sistema efectivamente integrado por el esclavo no es \eqref{eq:esclavo}, sino
su versión con la guía reconstruida:

\begin{equation}
\label{eq:esclavo-real}
\left\{
\begin{aligned}
  \dot{x}_s &= -\,y_s - z_s,\\
  \dot{y}_s &= x_s + a\,y_s + k\,\bigl(\,\hat{y}_m(t) - y_s\,\bigr),\\
  \dot{z}_s &= b + z_s\,(x_s - c).
\end{aligned}
\right.
\end{equation}

Esta distinción entre \eqref{eq:esclavo} y \eqref{eq:esclavo-real} es el eje de
todo el análisis que sigue.

\subsection{El spline cúbico}

La implementación construye $\hat{y}_m$ como el \emph{spline cúbico interpolante}
de los datos. Se recuerda su definición.

\begin{definicion}[Spline cúbico interpolante]
Dado $\mathcal{D}=\{(t_n,\eta_n)\}_{n=0}^{N-1}$ con $t_0<t_1<\cdots<t_{N-1}$, el
spline cúbico interpolante es la única función $S:[t_0,t_{N-1}]\to\R$ que satisface:
\begin{enumerate}
  \item[(i)] \textbf{Estructura por tramos:} en cada subintervalo
        $[t_n,t_{n+1}]$, $S$ coincide con un polinomio cúbico
        \[
          S_n(t)= \gamma_{0,n} + \gamma_{1,n}(t-t_n)
                 + \gamma_{2,n}(t-t_n)^2 + \gamma_{3,n}(t-t_n)^3 .
        \]
  \item[(ii)] \textbf{Interpolación:} $S(t_n)=\eta_n$ para todo $n$.
  \item[(iii)] \textbf{Regularidad:} $S \in C^{2}[t_0,t_{N-1}]$, es decir, $S$,
        $S'$ y $S''$ son continuas en los nodos interiores.
  \item[(iv)] \textbf{Condición de cierre:} dos condiciones adicionales en los
        extremos. La implementación emplea la condición \emph{not-a-knot}, que
        impone la continuidad de $S'''$ en $t_1$ y en $t_{N-2}$.
\end{enumerate}
\end{definicion}

El conteo de grados de libertad confirma que el problema está bien planteado:
$N-1$ tramos con $4$ coeficientes cada uno dan $4(N-1)$ incógnitas; las
condiciones (ii) aportan $2(N-1)$ ecuaciones, la continuidad de $S'$ y $S''$ en los
$N-2$ nodos interiores aporta $2(N-2)$, y las dos condiciones de cierre completan
$4(N-1)$. El sistema resultante es tridiagonal y de diagonal dominante, lo que
garantiza existencia, unicidad y estabilidad numérica de la solución.

Las tres propiedades de la Definición anterior son exactamente las que el problema
demanda:
\begin{itemize}
  \item la propiedad (i) permite evaluar $\hat{y}_m$ en cualquier instante, que era
        el requisito de la Sección \ref{sec:discreto};
  \item la propiedad (ii) asegura que en los nodos la reconstrucción es exacta, de
        modo que no se introduce sesgo sobre la información efectivamente recibida;
  \item la propiedad (iii) proporciona un miembro derecho de clase $C^2$ en $t$
        para la ecuación \eqref{eq:esclavo-real}, condición de regularidad que los
        métodos Runge--Kutta requieren para alcanzar su orden nominal.
\end{itemize}

\subsection{El error de reconstrucción}

\begin{definicion}[Error de reconstrucción]
Se define
\begin{equation}
\label{eq:delta}
  \delta(t) \;=\; \hat{y}_m(t) - y_m(t),
  \qquad t \in [0,T],
  \qquad
  \infnorm{\delta} \;=\; \sup_{t\in[0,T]} \lvert \delta(t)\rvert .
\end{equation}
\end{definicion}

Por la propiedad de interpolación, $\delta(t_n)=0$ en todos los nodos; el error
vive íntegramente en el interior de los subintervalos. La cantidad $\infnorm{\delta}$
es la única característica del interpolante que intervendrá en el resultado
principal.

% =====================================================================
\section{Resultado central: el error de sincronización está gobernado por $\delta$}
\label{sec:resultado}
% =====================================================================

\subsection{El sistema de error perturbado}

Se repite ahora el cálculo de la Sección \ref{sec:error}, pero partiendo del sistema
realmente integrado, \eqref{eq:esclavo-real}. La única ecuación que cambia es la de
$e_y$. Escribiendo $\hat{y}_m = y_m + \delta$:
\begin{align}
  \dot{e}_y
  &= e_x + a\,e_y + k\bigl(\hat{y}_m(t) - y_s\bigr) \nonumber\\
  &= e_x + a\,e_y + k\bigl(y_m(t) + \delta(t) - y_s\bigr) \nonumber\\
  &= e_x + (a-k)\,e_y \;+\; k\,\delta(t).
\end{align}

Las ecuaciones de $e_x$ y $e_z$ son idénticas a las anteriores. Por lo tanto:

\begin{equation}
\label{eq:error-perturbado}
\boxed{\;
\dot{e} \;=\; J(t)\,e \;+\; k\,\delta(t)\,u ,
\qquad
u = \begin{pmatrix} 0\\ 1\\ 0\end{pmatrix} .
\;}
\end{equation}

La lectura de \eqref{eq:error-perturbado} es directa y conceptualmente importante:
\emph{interpolar imperfectamente equivale a inyectar una perturbación externa en el
sistema de error}, con dos rasgos característicos. Primero, la perturbación actúa
exclusivamente en la dirección $u$ de la variable acoplada. Segundo, su amplitud es
$k\,\delta(t)$: la ganancia de acoplamiento, que estabiliza el sistema, amplifica
simultáneamente el error de reconstrucción.

\subsection{La cota del error de sincronización}

\begin{teorema}[Piso del error de sincronización]
\label{th:principal}
Supóngase que se cumple la condición de sincronización exponencial
\eqref{eq:cota-transicion} con constantes $M\ge 1$ y $\lambda>0$. Entonces la
solución de \eqref{eq:error-perturbado} satisface, para todo $t\ge 0$,
\begin{equation}
\label{eq:cota-completa}
  \dnorm{e(t)}
  \;\le\;
  \underbrace{M\,e^{-\lambda t}\,\dnorm{e(0)}}_{\text{(A) transitorio}}
  \;+\;
  \underbrace{\frac{k\,M}{\lambda}\,\bigl(1-e^{-\lambda t}\bigr)\,\infnorm{\delta}}_{\text{(B) piso}} ,
\end{equation}
y en consecuencia
\begin{equation}
\label{eq:cota-asintotica}
  \limsup_{t\to\infty}\;\dnorm{e(t)}
  \;\le\;
  \frac{k\,M}{\lambda}\;\infnorm{\delta}\, .
\end{equation}
\end{teorema}

\begin{proof}
La fórmula de variación de las constantes aplicada a
\eqref{eq:error-perturbado} da
\[
  e(t) \;=\; \Phi(t,0)\,e(0) \;+\; k\int_{0}^{t} \Phi(t,s)\,u\,\delta(s)\,\mathrm{d}s .
\]
Tomando normas y usando la desigualdad triangular junto con
$\dnorm{u}=1$ y $\lvert\delta(s)\rvert \le \infnorm{\delta}$:
\[
  \dnorm{e(t)}
  \;\le\; \dnorm{\Phi(t,0)}\,\dnorm{e(0)}
        + k\,\infnorm{\delta}\int_{0}^{t}\dnorm{\Phi(t,s)}\,\mathrm{d}s .
\]
Sustituyendo la cota \eqref{eq:cota-transicion} y evaluando la integral,
\[
  \int_{0}^{t} M\,e^{-\lambda(t-s)}\,\mathrm{d}s
  \;=\; \frac{M}{\lambda}\bigl(1-e^{-\lambda t}\bigr),
\]
se obtiene \eqref{eq:cota-completa}. Haciendo $t\to\infty$ el término (A) se anula
y $1-e^{-\lambda t}\to 1$, lo que produce \eqref{eq:cota-asintotica}.
\end{proof}

\subsection{Interpretación}

El Teorema \ref{th:principal} descompone el error de sincronización en dos
contribuciones de naturaleza completamente distinta:

\begin{itemize}
  \item El \textbf{término transitorio} (A) depende de la discrepancia de
        condiciones iniciales y decae exponencialmente. Se combate
        \emph{esperando}: es la función del periodo de sincronización
        $N_{\mathrm{sinc}}$. Ningún refinamiento del interpolante lo mejora, y
        tampoco hace falta.
  \item El \textbf{término de piso} (B) \emph{no decae}. Es una cota permanente,
        proporcional a $\infnorm{\delta}$, que persiste indefinidamente por mucho
        que se prolongue la integración. Aumentar $N_{\mathrm{sinc}}$ no lo reduce.
        La única palanca disponible sobre él es la calidad del operador de
        reconstrucción.
\end{itemize}

Esta separación responde con precisión a la pregunta que motiva el documento. La
interpolación es el mecanismo que actúa sobre el segundo término, es decir, sobre
el error residual que subsiste una vez agotado el transitorio. Cuanto menor sea
$\infnorm{\delta}$, más bajo será el piso al que puede descender el error de
sincronización.

Dos casos límite ilustran el alcance del resultado:
\[
  \delta \equiv 0
  \;\Longrightarrow\;
  \dnorm{e(t)} \le M e^{-\lambda t}\dnorm{e(0)} \to 0
  \quad\text{(sincronización idéntica)},
\]
mientras que con una reconstrucción burda el piso puede quedar por encima de la
tolerancia admisible para la aplicación, y en ese caso la sincronización
—aunque el criterio de Pecora y Carroll se cumpla— resulta inservible en la
práctica.

\begin{observacion}[El compromiso en la elección de $k$]
La ganancia $k$ aparece en el numerador de \eqref{eq:cota-asintotica} y, a la vez,
influye en $\lambda$ a través de la entrada $a-k$ de $J(t)$. Aumentar $k$ acelera
el transitorio (A) pero amplifica el piso (B). Este compromiso desaparece en la
medida en que $\infnorm{\delta}$ sea pequeño: con una reconstrucción de alto orden,
el término $k\,\infnorm{\delta}/\lambda$ permanece reducido incluso para ganancias
grandes, y $k$ puede elegirse atendiendo a la velocidad de convergencia sin pagar
un precio apreciable en precisión.
\end{observacion}

% =====================================================================
\section{El orden del interpolante y la magnitud del piso}
\label{sec:ordenes}
% =====================================================================

Establecido que el piso del error es proporcional a $\infnorm{\delta}$, resta
cuantificar $\infnorm{\delta}$ para distintas elecciones del operador
$\mathcal{I}$. Las cotas siguientes son resultados clásicos de teoría de la
aproximación; se enuncian sobre una malla uniforme de paso $h$ y para una función
$y_m$ suficientemente regular. Se emplea la notación
\[
  M_r \;=\; \infnorm{y_m^{(r)}} \;=\; \sup_{t\in[0,T]}\bigl\lvert y_m^{(r)}(t)\bigr\rvert .
\]

\subsection{Retención de orden cero}

La alternativa más simple consiste en congelar la última muestra disponible,
\[
  \hat{y}_m^{\,\mathrm{ZOH}}(t) = \eta_n \quad\text{para } t\in[t_n,t_{n+1}).
\]
Por el teorema del valor medio, para $t$ en dicho intervalo existe $\xi$ con
$\lvert y_m(t)-\eta_n\rvert = \lvert y_m'(\xi)\rvert\,(t-t_n)$, de donde
\begin{equation}
\label{eq:cota-zoh}
  \infnorm{\delta_{\mathrm{ZOH}}} \;\le\; M_1\,h \;=\; \mathcal{O}(h).
\end{equation}
Esta reconstrucción es discontinua en cada nodo. Además de su bajo orden, esa
discontinuidad viola la hipótesis de regularidad temporal del campo vectorial que
sostiene el análisis de orden de los métodos Runge--Kutta.

\subsection{Interpolación lineal por tramos}

Uniendo muestras consecutivas mediante segmentos rectos, el error clásico de
interpolación de Lagrange con dos nodos da
\begin{equation}
\label{eq:cota-lineal}
  \infnorm{\delta_{\mathrm{lin}}} \;\le\; \frac{M_2}{8}\,h^{2} \;=\; \mathcal{O}(h^{2}).
\end{equation}
El resultado es continuo, pero su derivada presenta saltos en los nodos: la
reconstrucción es $C^0$ y no $C^1$.

\subsection{Spline cúbico}

\begin{teorema}[Cota de error del spline cúbico]
\label{th:spline}
Sea $y_m \in C^{4}[t_0,t_{N-1}]$ y sea $S$ su spline cúbico interpolante sobre una
malla uniforme de paso $h$. Entonces
\begin{equation}
\label{eq:cota-spline}
  \infnorm{S - y_m} \;\le\; \frac{5}{384}\,M_4\,h^{4} \;=\; \mathcal{O}(h^{4}),
\end{equation}
y para la primera derivada
\begin{equation}
\label{eq:cota-spline-derivada}
  \infnorm{S' - y_m'} \;\le\; \frac{1}{24}\,M_4\,h^{3}.
\end{equation}
\end{teorema}

La hipótesis $y_m\in C^4$ se satisface holgadamente: como se señaló en la
Sección \ref{sec:planteamiento}, $y_m$ es solución de un sistema polinomial y por
tanto $y_m\in C^{\infty}$. Sobre el atractor —conjunto acotado— las derivadas de
todo orden permanecen acotadas, de modo que $M_4<\infty$. Es decir, la propia
naturaleza del sistema de Rössler es lo que legitima el empleo de un interpolante
de cuarto orden: la señal guía es exactamente el tipo de función para el que el
Teorema \ref{th:spline} rinde su máximo provecho.

\subsection{Comparación de las cotas}

Combinando el Teorema \ref{th:principal} con las cotas anteriores se obtiene, para
cada operador de reconstrucción, el piso correspondiente:

\begin{table}[H]
\centering
\caption{Regularidad, orden de aproximación y piso resultante del error de
sincronización según el operador de reconstrucción.}
\label{tab:ordenes}
\begin{tabular}{@{}llll@{}}
\toprule
\textbf{Reconstrucción} & \textbf{Regularidad} & $\infnorm{\delta}$ &
\textbf{Piso } $\displaystyle\limsup_t \dnorm{e(t)}$ \\
\midrule
Orden cero (retención) & discontinua & $M_1 h$              & $\dfrac{kM}{\lambda}M_1\,h$ \\[0.9em]
Lineal por tramos      & $C^{0}$     & $\dfrac{M_2}{8}h^{2}$ & $\dfrac{kM}{8\lambda}M_2\,h^{2}$ \\[0.9em]
\textbf{Spline cúbico} & $C^{2}$     & $\dfrac{5M_4}{384}h^{4}$ & $\dfrac{5kM}{384\lambda}M_4\,h^{4}$ \\
\bottomrule
\end{tabular}
\end{table}

El cociente entre el piso de la retención de orden cero y el del spline cúbico es
\begin{equation}
\label{eq:ganancia}
  \rho \;=\;
  \frac{M_1\,h}{\dfrac{5}{384}M_4\,h^{4}}
  \;=\;
  \frac{384}{5}\cdot\frac{M_1}{M_4}\cdot\frac{1}{h^{3}} .
\end{equation}

La estructura de \eqref{eq:ganancia} es lo relevante: el factor $h^{-3}$ crece sin
límite conforme se refina la malla. Con el paso empleado en el sistema, $h$ es del
orden de $10^{-2}$, de modo que solo el término $h^{-3}$ contribuye un factor de
$10^{6}$ a favor del spline. Debe subrayarse que \eqref{eq:ganancia} compara
\emph{cotas superiores} y no errores medidos; su valor es indicar el orden de
magnitud de la mejora estructural que aporta la interpolación de alto orden, no
predecir una cifra concreta.

\subsection{Consistencia con el orden del integrador}

Existe un segundo argumento, independiente del anterior, a favor del spline. El
método de integración empleado es de orden $5(4)$: su error local por paso se
comporta como $\mathcal{O}(\tau^{5})$ y su error global como
$\mathcal{O}(\tau^{4})$. Ahora bien, la precisión efectiva del resultado está
limitada por el peor de los dos términos que la componen:
\[
  \text{error total}
  \;\sim\;
  \underbrace{\mathcal{O}(\tau^{4})}_{\text{integrador}}
  \;+\;
  \underbrace{\mathcal{O}\bigl(\infnorm{\delta}\bigr)}_{\text{reconstrucción de la guía}} .
\]

Con una retención de orden cero, el segundo sumando es $\mathcal{O}(h)$ y domina
por completo: la precisión del conjunto degenera a primer orden y todo el esfuerzo
computacional invertido en un integrador de orden cinco se desperdicia. A ello se
añade que un miembro derecho discontinuo en $t$ invalida los desarrollos de Taylor
sobre los que se demuestran las condiciones de orden del método, por lo que ni
siquiera el orden reducido está garantizado de manera limpia.

Con el spline cúbico, en cambio, el segundo sumando es $\mathcal{O}(h^{4})$, del
mismo orden que el primero, y la reconstrucción de la guía deja de ser el eslabón
débil. La regularidad $C^{2}$ del spline provee además el miembro derecho suave
que la teoría del método presupone. Interpolación e integración quedan así
\emph{equilibradas en orden}, que es la condición de diseño correcta cuando se
acopla un reconstructor de señal a un integrador numérico.

% =====================================================================
\section{Consecuencia sobre el descifrado}
\label{sec:cripto}
% =====================================================================

Para cerrar el argumento conviene seguir el error de sincronización hasta su efecto
final sobre la imagen recuperada, ya que es allí donde se manifiesta la utilidad
práctica de todo lo anterior.

\subsection{Del error de sincronización al error de píxel}

Sea $d_j$ el $j$-ésimo valor del vector difundido en el maestro, $\ell_j$ el término
del mapa logístico y $\alpha$ el factor de escala. La etapa de confusión produce
\begin{equation}
\label{eq:cifrado}
  c_j \;=\; \alpha\,d_j \;+\; \ell_j \;+\; x_m(t_{n_j}),
\end{equation}
donde $t_{n_j}$ recorre los instantes posteriores al periodo de sincronización. El
esclavo regenera $\ell_j$ de forma exacta —el mapa logístico es una recurrencia
determinista que no depende de la sincronización— y despeja
\begin{equation}
\label{eq:descifrado}
  \hat{d}_j \;=\; \frac{1}{\alpha}\Bigl(\,c_j - \ell_j - x_s(t_{n_j})\,\Bigr).
\end{equation}

Restando \eqref{eq:descifrado} de la expresión exacta obtenida a partir de
\eqref{eq:cifrado}:
\begin{equation}
\label{eq:error-pixel}
\boxed{\;
  \hat{d}_j - d_j
  \;=\;
  \frac{1}{\alpha}\Bigl(x_m(t_{n_j}) - x_s(t_{n_j})\Bigr)
  \;=\;
  -\,\frac{1}{\alpha}\;e_x(t_{n_j}) .
\;}
\end{equation}

El error de reconstrucción de cada píxel es, exactamente, el error de
sincronización de la componente $x$ \emph{amplificado por el factor} $1/\alpha$.
Con el valor de escala empleado en el sistema, $\alpha = 10^{-2}$, dicho factor es
$100$: cualquier discrepancia en la sincronización se magnifica dos órdenes de
magnitud al llegar al dominio de la imagen.

\subsection{El umbral de cuantización}

La reconstrucción final cuantiza los valores a enteros de ocho bits mediante
$\mathrm{round}(255\,\hat{d}_j)$. El redondeo devuelve el entero correcto siempre
que el error se mantenga por debajo de medio nivel de cuantización:
\[
  255\,\bigl\lvert \hat{d}_j - d_j \bigr\rvert \;<\; \tfrac{1}{2}
  \qquad\Longleftrightarrow\qquad
  \bigl\lvert \hat{d}_j - d_j \bigr\rvert \;<\; \frac{1}{510}.
\]
Sustituyendo \eqref{eq:error-pixel} se obtiene la condición suficiente de
recuperación exacta:
\begin{equation}
\label{eq:umbral}
  \bigl\lvert e_x(t_{n_j}) \bigr\rvert \;<\; \frac{\alpha}{510}
  \;\approx\; 2\times 10^{-5}.
\end{equation}

Y combinando \eqref{eq:umbral} con el Teorema \ref{th:principal}, la exigencia
sobre el interpolante se vuelve explícita:
\begin{equation}
\label{eq:exigencia}
  \frac{k\,M}{\lambda}\,\infnorm{\delta} \;<\; \frac{\alpha}{510}.
\end{equation}

La desigualdad \eqref{eq:exigencia} resume el papel de la interpolación en el
sistema completo. El umbral del miembro derecho es estrecho, y el factor de
amplificación $1/\alpha$ lo estrecha aún más; un piso de error de orden
$\mathcal{O}(h)$ difícilmente lo satisface, mientras que un piso de orden
$\mathcal{O}(h^{4})$ deja un margen amplio. La interpolación por spline cúbico
es, en definitiva, lo que permite reducir el error de sincronización hasta situarlo
holgadamente por debajo del umbral que la cuantización impone, y con ello hacer
viable la recuperación fiel de la información cifrada.

% =====================================================================
\section{Detalles de implementación explicados por el análisis}
\label{sec:detalles}
% =====================================================================

Varias decisiones concretas del código encuentran su justificación en el desarrollo
anterior.

\paragraph{Mallas idénticas en ambos extremos.}
Maestro y esclavo construyen su malla temporal con la misma expresión, sobre
$N=N_{\mathrm{sinc}}+N_{\mathrm{key}}$ puntos equiespaciados en $[0,\,Nh]$. El paso
efectivo resultante es
\[
  h_{\mathrm{ef}} \;=\; \frac{N\,h}{N-1} \;=\; h\Bigl(1+\frac{1}{N-1}\Bigr),
\]
ligeramente superior al valor nominal $h$. Lo relevante no es el valor exacto sino
la coincidencia: al emplear ambos extremos la misma construcción, los nodos del
interpolante son los mismos instantes en los que el maestro evaluó la señal, y la
propiedad $\delta(t_n)=0$ se conserva. Una discrepancia entre las dos mallas
introduciría un desplazamiento sistemático en $\delta$ que ningún refinamiento del
interpolante podría corregir.

\paragraph{Extrapolación en la frontera.}
El interpolante se configura para admitir evaluación fuera de $[t_0,t_{N-1}]$. Esto
no responde a una necesidad de extender el dominio, sino a una particularidad del
paso adaptativo: al aproximarse al extremo final del intervalo, el integrador puede
proponer tentativamente un paso que rebase $t_{N-1}$ antes de reducirlo. Permitir
la extrapolación evita que dicha evaluación tentativa aborte la integración. El
número de puntos afectados es despreciable y ninguno de ellos contribuye al flujo
de clave.

\paragraph{Extracción del flujo de clave tras el transitorio.}
El vector de enmascaramiento se toma de la componente $x_s$ a partir del índice
$N_{\mathrm{sinc}}$. Según \eqref{eq:cota-completa}, ese índice delimita la región
temporal en la que el término (A) ya ha decaído y solo subsiste el piso (B). La
elección de $N_{\mathrm{sinc}}$ es, en esencia, la elección de un $t^{*}$ tal que
$M e^{-\lambda t^{*}}\dnorm{e(0)}$ resulte despreciable frente a
$k M \infnorm{\delta}/\lambda$.

\paragraph{Tolerancias del integrador.}
Maestro y esclavo emplean las mismas tolerancias relativa y absoluta en el control
de paso adaptativo. Es pertinente señalar que estas tolerancias establecen un piso
de error adicional, de origen distinto al analizado aquí. La interpolación elimina
la contribución correspondiente a la reconstrucción de la señal guía; una vez que
esa contribución deja de ser dominante, el error residual queda determinado por los
demás términos, y ajustar las tolerancias pasa a ser la palanca pertinente. El
análisis de este documento identifica cuál de los pisos controla la interpolación,
no afirma que sea el único existente.

% =====================================================================
\section{Síntesis}
% =====================================================================

El argumento completo puede resumirse en cinco proposiciones encadenadas.

\begin{enumerate}
  \item El esquema de Pecora y Carroll, en su forma acoplada \eqref{eq:esclavo},
        conduce a un sistema de error \emph{exactamente lineal}
        $\dot{e}=J(t)e$, que converge a cero de manera exponencial cuando los
        exponentes de Lyapunov condicionales del subsistema respuesta son
        negativos. Esta garantía presupone disponer de la señal guía como función
        continua del tiempo.

  \item El canal de comunicación entrega únicamente $N$ muestras discretas, en
        tanto que el integrador Runge--Kutta de paso adaptativo requiere evaluar la
        guía en instantes intermedios que no pertenecen a la malla. Sin un operador
        de reconstrucción, la ecuación del esclavo carece de definición en casi
        todo punto.

  \item La interpolación por spline cúbico construye ese operador: una función
        definida en todo el intervalo, exacta en los nodos y de clase $C^{2}$. Su
        discrepancia respecto de la señal verdadera es
        $\delta(t)=\hat{y}_m(t)-y_m(t)$.

  \item Al integrar con la guía reconstruida, el sistema de error se convierte en
        $\dot{e}=J(t)e+k\,\delta(t)\,u$, y su solución se descompone en un
        transitorio que decae exponencialmente —eliminado por el periodo de
        sincronización— y un piso permanente acotado por
        $k M \infnorm{\delta}/\lambda$. El piso es insensible al tiempo de
        integración y depende exclusivamente de la calidad de la interpolación.

  \item Las cotas clásicas de aproximación sitúan $\infnorm{\delta}$ en
        $\mathcal{O}(h)$ para una retención de orden cero y en $\mathcal{O}(h^{4})$
        para el spline cúbico, orden este último que además equilibra la
        reconstrucción con el orden del integrador. Aplicar la interpolación en el
        proceso de sincronización significa, por tanto, sustituir un piso de error
        de primer orden por uno de cuarto orden, lo que permite reducir el error de
        sincronización y situarlo por debajo del umbral que la cuantización de la
        imagen impone al descifrado.
\end{enumerate}

En una frase: dentro de este sistema, interpolar es reconstruir la señal guía que
el esquema de Pecora y Carroll exige y el canal no puede entregar, y el orden de
esa reconstrucción es lo que fija el error de sincronización residual.

\end{document}
